% Addendum to the Preface of BODYandSOUL: MATHEMATICAL SIMULATION TECHNOLOGY
% (draft (c) Claes Johnson 2010).  Written to follow the existing preface, using
% the series macros of bookdefs.tex.  Draft 2026.

\section*{Preface to the Reissue: the Reason for Change Has Arrived}

\begin{quote}
I admit that each and every thing remains in its state until there is reason for change.
\hfill (Leibniz)
\end{quote}

\noindent The preface above was written in 2010 and opens with that line. It then describes a
crisis with nobody responsible, a stalemate in which no move is possible, and names the cause: the
computer. Sixteen years later the diagnosis stands and the stalemate does not, because the reason
for change has arrived from outside the subject and does not require anybody in it to accept
responsibility.

\subsection*{What Has Happened}

In 2010 the argument of this book was that the computer had changed what mathematics is for, and
that education had not noticed. The remedy proposed was concrete and is stated in the preface: the
student should \emph{write the code} for the algorithm, and then let the machine do the tedious
routine work. Long division was the example --- do not drill the algorithm, implement it.

That remedy has itself been automated. Machines now write the code. They also carry out symbolic
derivation, and they are beginning to produce formal proofs that machine proof-checkers accept.
Every part of mathematics that consists of \emph{following a rule correctly} is now done better,
faster and more cheaply than a person can do it, and this includes the part this book asked the
student to take over from the pencil.

\subsection*{Why the Argument Survives}

Mathematics is two activities: \emph{formulate} the equation, and \emph{solve} it. Everything that
has been automated belongs to the second. Nothing has been automated in the first, and nothing
about the automation suggests it will be --- there is no rule for deciding which equation describes
a situation, because that decision is answerable to the world and not to a rule.

This is Leibniz's own distinction. Truths of \emph{reasoning} are those that follow from what is
already granted, and a machine that manipulates symbols correctly can produce them all. Truths of
\emph{fact} are not of that kind; they must be checked against how things are. What has been
automated is exactly the first and none of the second, which is why the question of what to teach
has an answer and not merely an anxiety.

\subsection*{The Quarrel of the Preface, Settled}

The preface above sets the constructive view --- $\sqrt2$ exists because its digits can be computed
--- against the formalist one, in which it exists because its non-existence would contradict an
axiom. That quarrel is now decided in an unexpected way. Formal derivation from axioms is precisely
what machines do well. It is not wrong; it has become \emph{cheap}, and cheapness settles curricula
where argument does not. What remains expensive, and therefore worth a student's years, is
constructing the model and judging whether the construction is right.

\subsection*{What Changes in This Book}

Less than one might fear. Parts II to V stand: a world described by differential equations, solved
by computation, is still what a student must understand, and understanding it is now the \emph{only}
reason to be in the room, since the typing is done elsewhere.

What changes is the student's relation to the machine. In 2010 the student wrote the simulator and
the machine ran it. Now the student \emph{specifies} the simulator, the machine writes it, and the
student must decide whether what comes back is correct. That is a harder task and not an easier
one, and it cannot be done by somebody who has only watched. The simulators in Parts III, VII and X
of this book are unusually well suited to it, because their answers are known: they can be
specified to a machine and the result checked against the analysis given here. Catching the machine
where it is wrong is the exercise, and it requires exactly the mathematics this book teaches.

\subsection*{On Responsibility}

The 2010 preface observes that no one is responsible and no one is permitted to become responsible.
That remains true and is now beside the point. A stalemate ends when the board changes, and the
board has changed. The question is no longer whether to reform mathematics education but whether
the reform will be designed or merely suffered.
